%%% Größen, Einheiten und Messen
\begin{lesen}
\lies{$l = 2{,}5\eh{m}$}{The length is two point five metres. Two point five is the numerical value, the metre is the unit.}{Die Länge beträgt zwei Komma fünf Meter. Zwei Komma fünf ist der Zahlenwert, Meter die Einheit.}
\lies{\skSkala}{Read the scale of the measuring instrument. Take the reading at eye level.}{Lies die Skala des Messgeräts ab. Lies den Messwert auf Augenhöhe ab.}
\lies{$t_1;\ t_2;\ \ldots;\ t_5$ \newline $\bar{t} = 1{,}42\eh{s}$}{Repeat the measurement five times and calculate the mean.}{Wiederhole die Messung fünfmal und berechne den Mittelwert.}
\lies{$F \sim s$}{The force is directly proportional to the extension: if the extension doubles, the force doubles.}{Die Kraft ist proportional zur Verlängerung: Verdoppelt sich die Verlängerung, so verdoppelt sich die Kraft.}
\lies{\skAusgleich}{Plot the readings and draw the line of best fit.}{Trage die Messwerte in ein Diagramm ein und zeichne die Ausgleichsgerade.}
\end{lesen}

%%% Bewegungen
\begin{lesen}
\lies{$s = 100\eh{m},\ t = 12{,}5\eh{s}$ \newline $v = \dfrac{s}{t} = 8\eh{m/s}$}{A sprinter covers one hundred metres in twelve point five seconds. Her average speed is eight metres per second.}{Eine Sprinterin legt hundert Meter in zwölf Komma fünf Sekunden zurück. Ihre Durchschnittsgeschwindigkeit beträgt acht Meter pro Sekunde.}
\lies{$a = \dfrac{20\eh{m/s}}{5\eh{s}} = 4\eh{m/s^2}$}{The car accelerates from rest to twenty metres per second in five seconds. Its acceleration is four metres per second squared.}{Das Auto beschleunigt aus der Ruhe in fünf Sekunden auf zwanzig Meter pro Sekunde. Seine Beschleunigung beträgt vier Meter pro Sekunde zum Quadrat.}
\lies{\skVTkonst}{The body moves at constant velocity: the velocity–time graph is a horizontal line.}{Der Körper bewegt sich gleichförmig: Das Geschwindigkeit-Zeit-Diagramm ist eine waagerechte Gerade.}
\lies{$s = \dfrac{g}{2} \cdot t^2 = 4{,}9\eh{m/s^2} \cdot (2\eh{s})^2$ \newline $s \approx 19{,}6\eh{m}$}{A stone is dropped from a bridge. After two seconds it has fallen about nineteen point six metres.}{Ein Stein wird von einer Brücke fallen gelassen. Nach zwei Sekunden ist er etwa neunzehn Komma sechs Meter gefallen.}
\lies{$T = 0{,}5\eh{s}$ \newline $f = \dfrac{1}{T} = 2\eh{Hz}$}{The wheel completes one revolution in half a second, so its frequency is two hertz.}{Das Rad braucht für eine Umdrehung eine halbe Sekunde, seine Frequenz beträgt also zwei Hertz.}
\end{lesen}

%%% Kräfte und ihre Wirkungen
\begin{lesen}
\lies{$m = 5\eh{kg}$ \newline $F_G = m \cdot g \approx 49\eh{N}$}{The weight of a five-kilogram mass is about forty-nine newtons.}{Die Gewichtskraft eines Körpers der Masse fünf Kilogramm beträgt etwa neunundvierzig Newton.}
\lies{$F = m \cdot a$ \newline $F = 1200\eh{kg} \cdot 2\eh{m/s^2} = 2400\eh{N}$}{A force of two thousand four hundred newtons is needed to accelerate the car at two metres per second squared.}{Um das Auto mit zwei Meter pro Sekunde zum Quadrat zu beschleunigen, ist eine Kraft von zweitausendvierhundert Newton nötig.}
\lies{$D = \dfrac{F}{s} = \dfrac{2\eh{N}}{0{,}05\eh{m}} = 40\eh{N/m}$}{A force of two newtons stretches the spring by five centimetres. The spring constant is forty newtons per metre.}{Eine Kraft von zwei Newton dehnt die Feder um fünf Zentimeter. Die Federkonstante beträgt vierzig Newton pro Meter.}
\lies{\skHebel}{The lever is balanced when the clockwise moment equals the anticlockwise moment.}{Der Hebel ist im Gleichgewicht, wenn rechtsdrehendes und linksdrehendes Drehmoment gleich groß sind.}
\lies{$p = \dfrac{F}{A} = \dfrac{600\eh{N}}{0{,}03\eh{m^2}}$ \newline $p = 20\,000\eh{Pa}$}{The pressure is the force divided by the area: twenty thousand pascals.}{Der Druck ist der Quotient aus Kraft und Fläche: zwanzigtausend Pascal.}
\lies{$\varrho = \dfrac{m}{V} = \dfrac{540\eh{g}}{200\eh{cm^3}}$ \newline $\varrho = 2{,}7\eh{g/cm^3}$}{The density of the block is two point seven grams per cubic centimetre: it is made of aluminium.}{Die Dichte des Quaders beträgt zwei Komma sieben Gramm pro Kubikzentimeter: Er besteht aus Aluminium.}
\end{lesen}

%%% Energie, Arbeit und Leistung
\begin{lesen}
\lies{$W = F \cdot s$ \newline $W = 50\eh{N} \cdot 4\eh{m} = 200\eh{J}$}{The work done is the force times the distance moved in the direction of the force: two hundred joules.}{Die verrichtete Arbeit ist Kraft mal Weg in Kraftrichtung: zweihundert Joule.}
\lies{$E_{\mathrm{pot}} = 2\eh{kg} \cdot 9{,}81\eh{m/s^2} \cdot 10\eh{m}$ \newline $E_{\mathrm{pot}} \approx 196\eh{J}$}{A two-kilogram mass is lifted ten metres. Its potential energy increases by about one hundred and ninety-six joules.}{Ein Körper der Masse zwei Kilogramm wird zehn Meter angehoben. Seine potentielle Energie nimmt um etwa hundertsechsundneunzig Joule zu.}
\lies{$E_{\mathrm{pot}} \rightarrow E_{\mathrm{kin}}$}{As the ball falls, potential energy is converted into kinetic energy.}{Beim Fallen des Balls wird potentielle Energie in kinetische Energie umgewandelt.}
\lies{$P = \dfrac{W}{t} = \dfrac{3000\eh{J}}{60\eh{s}} = 50\eh{W}$}{The power is the work done per unit time: fifty watts.}{Die Leistung ist die pro Zeit verrichtete Arbeit: fünfzig Watt.}
\lies{$\eta = \dfrac{30\eh{J}}{100\eh{J}} = 0{,}3 = 30\,\%$}{The efficiency of the engine is thirty per cent; seventy per cent of the energy is transferred as waste heat.}{Der Wirkungsgrad des Motors beträgt dreißig Prozent; siebzig Prozent der Energie werden als Abwärme abgegeben.}
\end{lesen}

%%% Wärmelehre
\begin{lesen}
\lies{$\vartheta = 20\eh{\degree C}$ \newline $T = 293{,}15\eh{K}$}{Twenty degrees Celsius corresponds to two hundred and ninety-three point one five kelvin.}{Zwanzig Grad Celsius entsprechen zweihundertdreiundneunzig Komma eins fünf Kelvin.}
\lies{$Q = c \cdot m \cdot \Delta T$ \newline $Q = 4{,}19\eh{\tfrac{kJ}{kg \cdot K}} \cdot 2\eh{kg} \cdot 30\eh{K}$ \newline $Q \approx 251\eh{kJ}$}{About two hundred and fifty-one kilojoules of heat are needed to raise the temperature of two kilograms of water by thirty kelvin.}{Um zwei Kilogramm Wasser um dreißig Kelvin zu erwärmen, sind etwa zweihunderteinundfünfzig Kilojoule Wärme nötig.}
\lies{\skFest \; \skFluessig \; \skGas}{In a solid the particles vibrate about fixed positions; in a liquid they can move past each other; in a gas they move freely.}{In einem festen Körper schwingen die Teilchen um feste Plätze; in einer Flüssigkeit sind sie gegeneinander verschiebbar; in einem Gas bewegen sie sich frei.}
\lies{$100\eh{\degree C}$}{Water boils at one hundred degrees Celsius. While it is boiling, its temperature stays constant.}{Wasser siedet bei hundert Grad Celsius. Während des Siedens bleibt die Temperatur konstant.}
\lies{Wärmeleitung}{Metals are good conductors of heat; air is a good insulator.}{Metalle sind gute Wärmeleiter; Luft ist ein guter Isolator.}
\end{lesen}

%%% Elektrizität und Magnetismus
\begin{lesen}
\lies{$R = \dfrac{U}{I} = \dfrac{12\eh{V}}{0{,}5\eh{A}} = 24\eh{\Omega}$}{The voltage across the resistor is twelve volts and the current through it is half an ampere. Its resistance is twenty-four ohms.}{Am Widerstand liegt eine Spannung von zwölf Volt an, die Stromstärke beträgt ein halbes Ampere. Der Widerstand beträgt vierundzwanzig Ohm.}
\lies{\skReihe}{The two resistors are connected in series: the current through them is the same.}{Die beiden Widerstände sind in Reihe geschaltet: Die Stromstärke ist überall gleich.}
\lies{\skParallel}{The two resistors are connected in parallel: the voltage across them is the same.}{Die beiden Widerstände sind parallel geschaltet: An beiden liegt die gleiche Spannung an.}
\lies{\sz{rmeter, t=A} \; \sz{rmeter, t=V}}{An ammeter is connected in series, a voltmeter in parallel with the component.}{Ein Strommesser wird in Reihe, ein Spannungsmesser parallel zum Bauteil geschaltet.}
\lies{\skMagnet}{Like poles repel each other, unlike poles attract each other.}{Gleichnamige Pole stoßen einander ab, ungleichnamige Pole ziehen einander an.}
\end{lesen}

%%% Schwingungen und Wellen
\begin{lesen}
\lies{$l = 1\eh{m}$ \newline $T = 2\pi \cdot \sqrt{\dfrac{l}{g}} \approx 2{,}0\eh{s}$}{A simple pendulum one metre long has a period of about two seconds.}{Ein Fadenpendel von einem Meter Länge hat eine Periodendauer von etwa zwei Sekunden.}
\lies{$f = 440\eh{Hz}$ \newline $\lambda = \dfrac{c}{f} = \dfrac{340\eh{m/s}}{440\eh{Hz}} \approx 0{,}77\eh{m}$}{The note A has a frequency of four hundred and forty hertz. In air its wavelength is about seventy-seven centimetres.}{Der Ton a hat eine Frequenz von vierhundertvierzig Hertz. In Luft beträgt seine Wellenlänge etwa siebenundsiebzig Zentimeter.}
\lies{\skAmplitude}{The larger the amplitude, the louder the sound; the higher the frequency, the higher the pitch.}{Je größer die Amplitude, desto lauter der Ton; je höher die Frequenz, desto höher der Ton.}
\lies{\skLaengswelle}{Sound is a longitudinal wave: the particles oscillate parallel to the direction of travel.}{Schall ist eine Längswelle: Die Teilchen schwingen parallel zur Ausbreitungsrichtung.}
\lies{$s = \dfrac{340\eh{m/s} \cdot 2\eh{s}}{2} = 340\eh{m}$}{The echo returns after two seconds, so the wall is three hundred and forty metres away.}{Das Echo kommt nach zwei Sekunden zurück, die Wand ist also dreihundertvierzig Meter entfernt.}
\end{lesen}

%%% Optik
\begin{lesen}
\lies{\skReflexion}{The angle of incidence equals the angle of reflection. Both angles are measured from the normal.}{Der Einfallswinkel ist gleich dem Reflexionswinkel. Beide Winkel werden zum Lot gemessen.}
\lies{\skBrechung}{When light passes from air into water, it is refracted towards the normal.}{Beim Übergang von Luft in Wasser wird das Licht zum Lot hin gebrochen.}
\lies{\skSammellinse}{Rays parallel to the principal axis are refracted through the focal point.}{Parallelstrahlen werden so gebrochen, dass sie durch den Brennpunkt verlaufen.}
\lies{$f = 10\eh{cm},\ g = 30\eh{cm}$ \newline $\dfrac{1}{b} = \dfrac{1}{f} - \dfrac{1}{g}$ \newline $b = 15\eh{cm}$}{The object is thirty centimetres from the lens. The image is formed fifteen centimetres behind the lens; it is real, inverted and diminished.}{Der Gegenstand steht dreißig Zentimeter vor der Linse. Das Bild entsteht fünfzehn Zentimeter hinter der Linse; es ist reell, umgekehrt und verkleinert.}
\lies{\skPrisma}{A prism splits white light into the colours of the spectrum.}{Ein Prisma zerlegt weißes Licht in die Spektralfarben.}
\end{lesen}

%%% Atom- und Kernphysik
\begin{lesen}
\lies{$^{14}_{\ 6}\mathrm{C}$}{carbon-fourteen: the nucleus contains six protons and eight neutrons}{Kohlenstoff-14: Der Kern enthält sechs Protonen und acht Neutronen}
\lies{$^{226}_{\ 88}\mathrm{Ra} \rightarrow {}^{222}_{\ 86}\mathrm{Rn} + {}^{4}_{2}\mathrm{He}$}{Radium-two two six decays into radon-two two two by emitting an alpha particle.}{Radium-226 zerfällt unter Aussendung eines Alphateilchens in Radon-222.}
\lies{$T_{1/2} = 8\eh{d}$}{The half-life of iodine-one three one is eight days: after eight days half of the nuclei have decayed, after sixteen days three quarters.}{Die Halbwertszeit von Iod-131 beträgt acht Tage: Nach acht Tagen ist die Hälfte der Kerne zerfallen, nach sechzehn Tagen drei Viertel.}
\lies{$\alpha,\ \beta,\ \gamma$}{Alpha radiation is stopped by a sheet of paper, beta radiation by a few millimetres of aluminium; gamma radiation is only reduced by thick lead.}{Alphastrahlung wird durch ein Blatt Papier abgeschirmt, Betastrahlung durch einige Millimeter Aluminium; Gammastrahlung wird erst durch dickes Blei geschwächt.}
\lies{$E = h \cdot f$}{The energy of a photon equals the Planck constant times the frequency.}{Die Energie eines Photons ist gleich dem planckschen Wirkungsquantum mal der Frequenz.}
\end{lesen}
